\documentclass[11pt]{article}
\usepackage[margin=0.9in]{geometry}
\usepackage{amsmath, amssymb}
\usepackage{enumitem}
\usepackage{tcolorbox}
\usepackage{array}
\usepackage{longtable}
\usepackage{booktabs}
\usepackage{multicol}

\setlength{\parindent}{0pt}
\setlength{\parskip}{4pt}

\title{\textbf{CAT Quant Modern Math Practice Bank: P\&C and Probability}}
\author{}
\date{}

\begin{document}
\maketitle

\section*{Instructions}

\begin{itemize}
\item This is an original CAT-level Modern Math practice bank focused on P\&C, Counting, and Probability.
\item It is not an official, leaked, or guaranteed CAT paper.
\item Suggested time: 160 to 190 minutes for the full 20-question bank.
\item Calculators are not assumed.
\item Attempt under timed conditions before checking solutions.
\item Questions 1--10 cover P\&C / Counting / Arrangement / Selection. Questions 11--20 cover Probability.
\item For TITA questions, enter the exact numerical value. For MCQ, exactly one option is correct.
\end{itemize}

\newpage

\section*{Question Paper}

\subsection*{Question 1}
\textbf{Topic:} P\&C / Counting \\
\textbf{Subtopic:} Arrangement with repeated objects \\
\textbf{Difficulty:} Medium-Hard \\
\textbf{Type:} MCQ \\
\textbf{Estimated Time:} 3--4 minutes

\begin{tcolorbox}
All the letters of the word \textbf{MISSION} are arranged in a row using each letter exactly once. In how many such arrangements is it true that the two S's are not next to each other and the two I's are also not next to each other?
\end{tcolorbox}

\begin{enumerate}[label=(\Alph*)]
\item 540
\item 600
\item 660
\item 720
\end{enumerate}

\subsection*{Question 2}
\textbf{Topic:} P\&C / Counting \\
\textbf{Subtopic:} Distribution of identical objects with upper and lower bounds \\
\textbf{Difficulty:} Hard \\
\textbf{Type:} TITA \\
\textbf{Estimated Time:} 4--5 minutes

\begin{tcolorbox}
Twenty identical chocolates are to be distributed among four children A, B, C, and D such that each child receives at least 2 chocolates and at most 8 chocolates. Find the number of ways this distribution can be done.
\end{tcolorbox}

\textbf{Type:} TITA

\subsection*{Question 3}
\textbf{Topic:} P\&C / Counting \\
\textbf{Subtopic:} Circular arrangement with restriction \\
\textbf{Difficulty:} Hard \\
\textbf{Type:} MCQ \\
\textbf{Estimated Time:} 4--5 minutes

\begin{tcolorbox}
Eight people, including three friends X, Y, and Z, are to be seated around a circular table. In how many ways can they be seated such that no two of X, Y, and Z occupy adjacent seats?
\end{tcolorbox}

\begin{enumerate}[label=(\Alph*)]
\item 1200
\item 1440
\item 2160
\item 2520
\end{enumerate}

\subsection*{Question 4}
\textbf{Topic:} P\&C / Counting \\
\textbf{Subtopic:} Counting integer solutions under constraints \\
\textbf{Difficulty:} Hard \\
\textbf{Type:} MCQ \\
\textbf{Estimated Time:} 4--5 minutes

\begin{tcolorbox}
The number of ordered triples \( (a, b, c) \) of positive integers satisfying \( a < b < c \) and \( a + b + c = 30 \) is:
\end{tcolorbox}

\begin{enumerate}[label=(\Alph*)]
\item 54
\item 61
\item 67
\item 75
\end{enumerate}

\subsection*{Question 5}
\textbf{Topic:} P\&C / Counting \\
\textbf{Subtopic:} Committee formation with representation constraints \\
\textbf{Difficulty:} Hard \\
\textbf{Type:} MCQ \\
\textbf{Estimated Time:} 4--5 minutes

\begin{tcolorbox}
From 6 men and 5 women, a committee of 6 members is to be formed such that the committee contains at least 2 men and at least 2 women. From the formed committee, one member is to be chosen as the Chairperson and a different member as the Secretary. In how many ways can the committee be formed and the two offices be filled?
\end{tcolorbox}

\begin{enumerate}[label=(\Alph*)]
\item 11900
\item 12500
\item 12750
\item 13200
\end{enumerate}

\subsection*{Question 6}
\textbf{Topic:} P\&C / Counting \\
\textbf{Subtopic:} Digit/string counting with position restrictions \\
\textbf{Difficulty:} Hard \\
\textbf{Type:} TITA \\
\textbf{Estimated Time:} 3--4 minutes

\begin{tcolorbox}
How many 5-digit positive integers are there in which the digits are strictly increasing from left to right and which contain the digit 5?
\end{tcolorbox}

\textbf{Type:} TITA

\subsection*{Question 7}
\textbf{Topic:} P\&C / Counting \\
\textbf{Subtopic:} Distribution of identical objects with mixed parity and bound restrictions \\
\textbf{Difficulty:} Very Hard \\
\textbf{Type:} MCQ \\
\textbf{Estimated Time:} 6--7 minutes

\begin{tcolorbox}
Twelve identical balls are to be placed into four distinct boxes \( B_1, B_2, B_3, B_4 \), where every ball must be placed in some box. The placement must satisfy all the following conditions simultaneously:
\begin{itemize}
\item \( B_1 \) contains at most 3 balls,
\item \( B_2 \) contains at least 2 balls,
\item \( B_3 \) contains an odd number of balls,
\item \( B_4 \) contains an even number of balls (zero is allowed).
\end{itemize}
The number of valid distributions is:
\end{tcolorbox}

\begin{enumerate}[label=(\Alph*)]
\item 42
\item 48
\item 50
\item 56
\end{enumerate}

\subsection*{Question 8}
\textbf{Topic:} P\&C / Counting \\
\textbf{Subtopic:} Casework with leading-digit restriction \\
\textbf{Difficulty:} Hard \\
\textbf{Type:} TITA \\
\textbf{Estimated Time:} 4--5 minutes

\begin{tcolorbox}
How many 4-digit positive integers contain the digit 3 at least twice (counting positions)?
\end{tcolorbox}

\textbf{Type:} TITA

\subsection*{Question 9}
\textbf{Topic:} P\&C / Counting \\
\textbf{Subtopic:} Linear arrangement with position restrictions \\
\textbf{Difficulty:} Medium-Hard \\
\textbf{Type:} MCQ \\
\textbf{Estimated Time:} 3--4 minutes

\begin{tcolorbox}
Eight people, including two specific persons A and B, are to be seated in a row of eight chairs. In how many such arrangements are A and B both seated at non-end positions and also not seated next to each other?
\end{tcolorbox}

\begin{enumerate}[label=(\Alph*)]
\item 12960
\item 13800
\item 14400
\item 15120
\end{enumerate}

\subsection*{Question 10}
\textbf{Topic:} P\&C / Counting \\
\textbf{Subtopic:} Counting after removing invalid cases (multiset structure) \\
\textbf{Difficulty:} Very Hard \\
\textbf{Type:} TITA \\
\textbf{Estimated Time:} 5--6 minutes

\begin{tcolorbox}
A 4-letter ``word'' (an ordered sequence of 4 letters) is to be formed by choosing each letter from the set \( \{A, B, C, D, E\} \) with repetition allowed. How many such words have the property that exactly one letter of the alphabet appears two or more times in the word, while every other letter that appears in the word appears at most once?
\end{tcolorbox}

\textbf{Type:} TITA

\subsection*{Question 11}
\textbf{Topic:} Probability \\
\textbf{Subtopic:} Conditional probability \\
\textbf{Difficulty:} Medium-Hard \\
\textbf{Type:} MCQ \\
\textbf{Estimated Time:} 3--4 minutes

\begin{tcolorbox}
A bag contains 5 red balls and 4 blue balls. Three balls are drawn at random without replacement. Given that at least one of the drawn balls is blue, what is the probability that exactly two of the drawn balls are blue?
\end{tcolorbox}

\begin{enumerate}[label=(\Alph*)]
\item \( \dfrac{5}{14} \)
\item \( \dfrac{15}{37} \)
\item \( \dfrac{6}{19} \)
\item \( \dfrac{9}{28} \)
\end{enumerate}

\subsection*{Question 12}
\textbf{Topic:} Probability \\
\textbf{Subtopic:} Probability after careful counting (parity-based selection) \\
\textbf{Difficulty:} Hard \\
\textbf{Type:} MCQ \\
\textbf{Estimated Time:} 4--5 minutes

\begin{tcolorbox}
Four distinct numbers are chosen at random from the set \( \{1, 2, 3, \ldots, 12\} \). What is the probability that the sum of the four chosen numbers is even?
\end{tcolorbox}

\begin{enumerate}[label=(\Alph*)]
\item \( \dfrac{1}{2} \)
\item \( \dfrac{5}{11} \)
\item \( \dfrac{16}{33} \)
\item \( \dfrac{17}{33} \)
\end{enumerate}

\subsection*{Question 13}
\textbf{Topic:} Probability \\
\textbf{Subtopic:} Conditional probability with restricted sample space \\
\textbf{Difficulty:} Hard \\
\textbf{Type:} MCQ \\
\textbf{Estimated Time:} 5--6 minutes

\begin{tcolorbox}
Three distinct numbers are chosen at random (without replacement) from the set \( \{1, 2, 3, \ldots, 12\} \). Given that the smallest of the three chosen numbers is at most 4, what is the probability that the three chosen numbers form an arithmetic progression (in some order)?
\end{tcolorbox}

\begin{enumerate}[label=(\Alph*)]
\item \( \dfrac{1}{10} \)
\item \( \dfrac{9}{110} \)
\item \( \dfrac{9}{82} \)
\item \( \dfrac{18}{55} \)
\end{enumerate}

\subsection*{Question 14}
\textbf{Topic:} Probability \\
\textbf{Subtopic:} Selection-based probability with multiset structure \\
\textbf{Difficulty:} Hard \\
\textbf{Type:} TITA \\
\textbf{Estimated Time:} 4--5 minutes

\begin{tcolorbox}
Three letters are picked at random one by one, without replacement, from the eight letters of the word \textbf{PARALLEL}. What is the probability that the three picked letters are all different letters of the alphabet? (Express your answer as a simplified fraction in the form \( p/q \).)
\end{tcolorbox}

\textbf{Type:} TITA

\subsection*{Question 15}
\textbf{Topic:} Probability \\
\textbf{Subtopic:} Bayes-style conditional probability over a mixed experiment \\
\textbf{Difficulty:} Hard \\
\textbf{Type:} MCQ \\
\textbf{Estimated Time:} 5--6 minutes

\begin{tcolorbox}
A fair coin is to be tossed \( n \) times, where the value of \( n \) is chosen uniformly at random from \( \{3, 4, 5\} \) before the experiment starts. After the experiment, you are told that exactly 2 of the tosses came up heads. What is the probability that the chosen value of \( n \) was 4?
\end{tcolorbox}

\begin{enumerate}[label=(\Alph*)]
\item \( \dfrac{1}{3} \)
\item \( \dfrac{5}{17} \)
\item \( \dfrac{6}{17} \)
\item \( \dfrac{6}{19} \)
\end{enumerate}

\subsection*{Question 16}
\textbf{Topic:} Probability \\
\textbf{Subtopic:} Probability from random ordering (dice with positional constraint) \\
\textbf{Difficulty:} Hard \\
\textbf{Type:} TITA \\
\textbf{Estimated Time:} 4--5 minutes

\begin{tcolorbox}
A fair six-sided die is rolled three times. What is the probability that the number obtained on the second roll is strictly greater than the numbers obtained on both the first roll and the third roll? (Express your answer as a simplified fraction in the form \( p/q \).)
\end{tcolorbox}

\textbf{Type:} TITA

\subsection*{Question 17}
\textbf{Topic:} Probability \\
\textbf{Subtopic:} Probability with dependent events / waiting-time setup \\
\textbf{Difficulty:} Very Hard \\
\textbf{Type:} MCQ \\
\textbf{Estimated Time:} 6--7 minutes

\begin{tcolorbox}
A bag contains 6 red balls and 4 blue balls. Balls are drawn one by one, without replacement, until a red ball is drawn for the second time. The drawing then stops. What is the probability that the drawing stops on exactly the 5th draw?
\end{tcolorbox}

\begin{enumerate}[label=(\Alph*)]
\item \( \dfrac{1}{21} \)
\item \( \dfrac{2}{21} \)
\item \( \dfrac{1}{7} \)
\item \( \dfrac{4}{21} \)
\end{enumerate}

\subsection*{Question 18}
\textbf{Topic:} Probability \\
\textbf{Subtopic:} Arrangement-based probability with adjacency restrictions \\
\textbf{Difficulty:} Very Hard \\
\textbf{Type:} TITA \\
\textbf{Estimated Time:} 6--7 minutes

\begin{tcolorbox}
Five cards numbered 1, 2, 3, 4, and 5 are shuffled and placed in a row in a uniformly random order. What is the probability that, in the resulting arrangement, no two cards whose numbers differ by exactly 1 are placed next to each other? (Express your answer as a simplified fraction in the form \( p/q \).)
\end{tcolorbox}

\textbf{Type:} TITA

\subsection*{Question 19}
\textbf{Topic:} Probability \\
\textbf{Subtopic:} Probability with divisibility (inclusion-exclusion based) \\
\textbf{Difficulty:} Medium-Hard \\
\textbf{Type:} MCQ \\
\textbf{Estimated Time:} 3--4 minutes

\begin{tcolorbox}
A number is chosen uniformly at random from the set \( \{1, 2, 3, \ldots, 200\} \). What is the probability that the chosen number is divisible by exactly two of the three numbers 2, 3, and 5?
\end{tcolorbox}

\begin{enumerate}[label=(\Alph*)]
\item \( \dfrac{6}{25} \)
\item \( \dfrac{7}{25} \)
\item \( \dfrac{11}{50} \)
\item \( \dfrac{13}{50} \)
\end{enumerate}

\subsection*{Question 20}
\textbf{Topic:} Probability \\
\textbf{Subtopic:} Probability combining counting and ordering (with replacement) \\
\textbf{Difficulty:} Hard \\
\textbf{Type:} TITA \\
\textbf{Estimated Time:} 4--5 minutes

\begin{tcolorbox}
A box contains 5 balls labelled 1, 2, 3, 4, and 5. A ball is drawn at random, its label is noted, and it is then returned to the box. This procedure is performed three times in total. What is the probability that the three noted labels are all distinct AND form an arithmetic progression in some order? (Express your answer as a simplified fraction in the form \( p/q \).)
\end{tcolorbox}

\textbf{Type:} TITA

\newpage

\section*{Answer Key}

\begin{longtable}{|c|l|l|c|c|c|}
\hline
\textbf{Q\#} & \textbf{Topic} & \textbf{Subtopic} & \textbf{Type} & \textbf{Answer} & \textbf{Difficulty} \\
\hline
1 & P\&C & Arrangement (repeated objects) & MCQ & (C) 660 & Medium-Hard \\
\hline
2 & P\&C & Distribution with bounds & TITA & 231 & Hard \\
\hline
3 & P\&C & Circular arrangement & MCQ & (B) 1440 & Hard \\
\hline
4 & P\&C & Integer solutions & MCQ & (B) 61 & Hard \\
\hline
5 & P\&C & Committee with offices & MCQ & (C) 12750 & Hard \\
\hline
6 & P\&C & Digit counting & TITA & 70 & Hard \\
\hline
7 & P\&C & Distribution (parity + bounds) & MCQ & (C) 50 & Very Hard \\
\hline
8 & P\&C & Casework / 4-digit counting & TITA & 495 & Hard \\
\hline
9 & P\&C & Linear arrangement (restrictions) & MCQ & (C) 14400 & Medium-Hard \\
\hline
10 & P\&C & Multiset-based counting & TITA & 445 & Very Hard \\
\hline
11 & Probability & Conditional probability & MCQ & (B) 15/37 & Medium-Hard \\
\hline
12 & Probability & Parity-based selection & MCQ & (D) 17/33 & Hard \\
\hline
13 & Probability & Conditional (restricted space) & MCQ & (C) 9/82 & Hard \\
\hline
14 & Probability & Multiset selection & TITA & 17/28 & Hard \\
\hline
15 & Probability & Bayes (mixed experiment) & MCQ & (C) 6/17 & Hard \\
\hline
16 & Probability & Dice (positional) & TITA & 55/216 & Hard \\
\hline
17 & Probability & Waiting time / dependent events & MCQ & (B) 2/21 & Very Hard \\
\hline
18 & Probability & Arrangement (no adjacent) & TITA & 7/60 & Very Hard \\
\hline
19 & Probability & Divisibility-based probability & MCQ & (A) 6/25 & Medium-Hard \\
\hline
20 & Probability & Counting + ordering (with rep.) & TITA & 24/125 & Hard \\
\hline
\end{longtable}

\newpage

\section*{Detailed Solutions}

\subsection*{Solution 1}
\textbf{Answer:} (C) 660

\textbf{Detailed Solution:}

Step 1: The word MISSION has 7 letters: M, I, S, S, I, O, N --- two S's, two I's, and three distinct letters M, O, N.

Step 2: Total arrangements of all letters \( = \dfrac{7!}{2! \cdot 2!} = \dfrac{5040}{4} = 1260 \).

Step 3: Let \( P \) = arrangements where the two S's are adjacent. Treat SS as one block: units are \([SS], I, I, M, O, N\) (6 units with two I's). So \( |P| = \dfrac{6!}{2!} = 360 \).

Step 4: By symmetry, let \( Q \) = arrangements with two I's adjacent. \( |Q| = 360 \).

Step 5: Let \( P \cap Q \) = arrangements with both pairs adjacent. Units: \([SS], [II], M, O, N\) (5 distinct units). \( |P \cap Q| = 5! = 120 \).

Step 6: By inclusion-exclusion, arrangements with at least one pair adjacent \( = 360 + 360 - 120 = 600 \). Required arrangements \( = 1260 - 600 = 660 \).

\textbf{Why this is CAT-level:} Two interacting ``no two adjacent'' restrictions push the student towards the gap method, which becomes messy because the gap method has to handle both letter pairs simultaneously. Complement-via-inclusion-exclusion is cleaner.

\textbf{Common Wrong Approach:} Subtracting only one ``together'' case and forgetting the double-overlap correction, or using gap method incorrectly and dropping cases.

\subsection*{Solution 2}
\textbf{Answer:} 231

\textbf{Detailed Solution:}

Step 1: Let \( a, b, c, d \) denote the chocolates received by A, B, C, D. We need \( a + b + c + d = 20 \) with \( 2 \le a, b, c, d \le 8 \).

Step 2: Substitute \( a = a' + 2 \), etc. Then \( a' + b' + c' + d' = 12 \) with \( 0 \le a', b', c', d' \le 6 \).

Step 3: Without the upper bound, the count is \( \binom{12 + 3}{3} = \binom{15}{3} = 455 \).

Step 4: Subtract cases where some variable exceeds 6 (i.e., \( a' \ge 7 \)). Substituting \( a'' = a' - 7 \), we need \( a'' + b' + c' + d' = 5 \), giving \( \binom{8}{3} = 56 \) solutions. There are 4 such variables, so subtract \( 4 \times 56 = 224 \).

Step 5: Two variables \( \ge 7 \) would require sum \( \ge 14 > 12 \), impossible. So no further correction.

Step 6: Total \( = 455 - 224 = 231 \).

\textbf{Why this is CAT-level:} The upper-bound condition makes raw stars-and-bars wrong; the trap is to skip inclusion-exclusion or apply it incorrectly.

\textbf{Common Wrong Approach:} Reporting 455 without subtracting the upper-bound violations, or double-subtracting.

\subsection*{Solution 3}
\textbf{Answer:} (B) 1440

\textbf{Detailed Solution:}

Step 1: First seat the 5 people who are not in \( \{X, Y, Z\} \) around the circle. The number of distinct circular arrangements is \( (5-1)! = 24 \).

Step 2: These 5 people create exactly 5 gaps (between consecutive seated people) into which X, Y, Z can be inserted.

Step 3: To ensure that no two of X, Y, Z are adjacent, each of X, Y, Z must be placed in a different gap (placing two of them in the same gap would make them adjacent).

Step 4: Number of ways to choose 3 distinct gaps out of 5 and arrange X, Y, Z in them \( = P(5,3) = 5 \times 4 \times 3 = 60 \).

Step 5: Required count \( = 24 \times 60 = 1440 \).

\textbf{Why this is CAT-level:} A naive student might try ``total \( - \) at least two adjacent'' and get tangled in overlapping events. The gap method gives a clean answer.

\textbf{Common Wrong Approach:} Using \( (8-1)! - \) (cases where X, Y, Z all together) and forgetting other adjacency patterns.

\subsection*{Solution 4}
\textbf{Answer:} (B) 61

\textbf{Detailed Solution:}

Step 1: We need ordered triples \( (a, b, c) \) with \( a < b < c \), \( a + b + c = 30 \), all positive integers. ``Ordered'' under \( a<b<c \) means each unordered triple gives exactly one such ordered triple.

Step 2: For each value of \( a \), we count pairs \( (b, c) \) with \( a < b < c \) and \( b + c = 30 - a \).

Step 3: Given \( a \), the smallest value of \( b \) is \( a+1 \) and we need \( b < c = (30-a) - b \), i.e., \( b < (30 - a)/2 \). Also \( b \ge a+1 \). The count is the number of integers \( b \) with \( a + 1 \le b < (30 - a)/2 \).

Step 4: Since \( c > b \) forces \( b \le \lfloor (29 - a)/2 \rfloor \). Enumerate:

\( a=1: b \in \{2,\ldots,14\} \) \( \Rightarrow 13 \) values. \\
\( a=2: b \in \{3,\ldots,13\} \) \( \Rightarrow 11 \). \\
\( a=3: \{4,\ldots,13\} \Rightarrow 10 \). \\
\( a=4: \{5,\ldots,12\} \Rightarrow 8 \). \\
\( a=5: \{6,\ldots,12\} \Rightarrow 7 \). \\
\( a=6: \{7,\ldots,11\} \Rightarrow 5 \). \\
\( a=7: \{8,\ldots,11\} \Rightarrow 4 \). \\
\( a=8: \{9,10\} \Rightarrow 2 \). \\
\( a=9: \{10\} \Rightarrow 1 \). \\
For \( a \ge 10 \), \( a < b < c \) forces \( a+b+c > 30 \).

Step 5: Total \( = 13+11+10+8+7+5+4+2+1 = 61 \).

\textbf{Why this is CAT-level:} The student may try a substitution approach and lose track of bounds. Direct casework on \( a \) is clean but requires careful bookkeeping of the bound on \( b \).

\textbf{Common Wrong Approach:} Counting positive integer solutions of \( a + b + c = 30 \) and dividing by \( 3! \), which incorrectly handles triples with repeated values.

\subsection*{Solution 5}
\textbf{Answer:} (C) 12750

\textbf{Detailed Solution:}

Step 1: Let \( (m, w) \) denote the number of men and women on the committee with \( m + w = 6 \), \( m \ge 2 \), \( w \ge 2 \).

Step 2: The valid splits are:
\( (2, 4): \binom{6}{2}\binom{5}{4} = 15 \cdot 5 = 75. \) \\
\( (3, 3): \binom{6}{3}\binom{5}{3} = 20 \cdot 10 = 200. \) \\
\( (4, 2): \binom{6}{4}\binom{5}{2} = 15 \cdot 10 = 150. \)

Step 3: Total committees \( = 75 + 200 + 150 = 425 \).

Step 4: For each committee of 6, choose a Chairperson (6 ways) and a different Secretary (5 ways): \( 6 \times 5 = 30 \).

Step 5: Total ways \( = 425 \times 30 = 12750 \).

\textbf{Why this is CAT-level:} The student must remember to exclude \( (5,1) \) and \( (6,0) \) cases (and the symmetric (1,5)) and then layer in the ordered office selection.

\textbf{Common Wrong Approach:} Counting unordered office pairs (\( \binom{6}{2} \) instead of \( 6 \times 5 \)) or including illegal compositions.

\subsection*{Solution 6}
\textbf{Answer:} 70

\textbf{Detailed Solution:}

Step 1: A 5-digit number with strictly increasing digits cannot start with 0 (otherwise 0 must precede all other digits, contradicting strict increase). So all digits are from \( \{1, 2, \ldots, 9\} \).

Step 2: Any subset of 5 digits from \( \{1,\ldots,9\} \) gives exactly one strictly increasing arrangement. So total strictly increasing 5-digit numbers \( = \binom{9}{5} = 126 \).

Step 3: Strictly increasing 5-digit numbers without the digit 5: choose 5 digits from \( \{1,2,3,4,6,7,8,9\} \), i.e., \( \binom{8}{5} = 56 \).

Step 4: Required count \( = 126 - 56 = 70 \).

\textbf{Why this is CAT-level:} The student must recognise that strict monotonicity collapses arrangement to selection, and must use complement to handle ``contains 5''.

\textbf{Common Wrong Approach:} Trying to directly place 5 in each of 5 positions and overcounting, or including 0 as a possible digit.

\subsection*{Solution 7}
\textbf{Answer:} (C) 50

\textbf{Detailed Solution:}

Step 1: Let \( a, b, c, d \) be the contents of \( B_1, B_2, B_3, B_4 \) with \( a + b + c + d = 12 \), \( 0 \le a \le 3 \), \( b \ge 2 \), \( c \) odd \( \ge 1 \), \( d \) even \( \ge 0 \).

Step 2: Substitute \( b' = b - 2 \), \( c = 2k+1 \) (\( k \ge 0 \)), \( d = 2m \) (\( m \ge 0 \)). Then \( a + b' + 2(k+m) = 9 \), with \( 0 \le a \le 3 \), \( b' \ge 0 \), \( k, m \ge 0 \).

Step 3: Let \( s = k + m \). For each \( s \), the number of \( (k, m) \) pairs with \( k + m = s \) is \( s + 1 \).

Step 4: For fixed \( (a, s) \): \( b' = 9 - 2s - a \), requiring \( 9 - 2s - a \ge 0 \), i.e., \( a \le 9 - 2s \). Combined with \( a \le 3 \):

\( s=0: a \le 9 \Rightarrow a \in \{0,1,2,3\} \), 4 values; \( (k,m) \) pairs: 1. Contribution: \( 4 \). \\
\( s=1: a \le 7 \Rightarrow 4 \) values; 2 pairs. Contribution: \( 8 \). \\
\( s=2: a \le 5 \Rightarrow 4 \) values; 3 pairs. Contribution: \( 12 \). \\
\( s=3: a \le 3 \Rightarrow 4 \) values; 4 pairs. Contribution: \( 16 \). \\
\( s=4: a \le 1 \Rightarrow a \in \{0,1\} \), 2 values; 5 pairs. Contribution: \( 10 \). \\
For \( s \ge 5 \), \( 9 - 2s < 0 \), no valid \( a \).

Step 5: Total \( = 4 + 8 + 12 + 16 + 10 = 50 \).

\textbf{Why this is CAT-level:} Four interacting conditions of different types (upper bound, lower bound, parity, parity) force the student to combine parity substitution with bounded stars-and-bars. Direct formula does not apply.

\textbf{Common Wrong Approach:} Ignoring the upper bound on \( B_1 \) or treating odd/even constraints loosely.

\subsection*{Solution 8}
\textbf{Answer:} 495

\textbf{Detailed Solution:}

Step 1: A 4-digit number has digits \( d_1 d_2 d_3 d_4 \) with \( d_1 \in \{1, \ldots, 9\} \) and others in \( \{0, \ldots, 9\} \). We count those with at least two 3's by counting exactly 2, exactly 3, exactly 4.

Step 2 (exactly 2 threes):
\begin{itemize}
\item Case A: \( d_1 = 3 \). Choose 1 of the remaining 3 positions to also be 3: \( \binom{3}{1}=3 \). The other two positions are non-3 digits: \( 9 \times 9 = 81 \) ways. Subtotal: \( 3 \times 81 = 243 \).
\item Case B: \( d_1 \neq 3 \). Then \( d_1 \in \{1,2,4,5,6,7,8,9\} \) (8 choices). Two of \( d_2, d_3, d_4 \) are 3 (\( \binom{3}{2}=3 \) ways) and the remaining position is a non-3 digit (9 ways). Subtotal: \( 8 \times 3 \times 9 = 216 \).
\end{itemize}
Subtotal: \( 243 + 216 = 459 \).

Step 3 (exactly 3 threes):
\begin{itemize}
\item Case A: \( d_1 = 3 \). Choose 2 of the other 3 positions to be 3 (\( \binom{3}{2} = 3 \)). Remaining position is non-3 digit (9 choices). Subtotal: \( 3 \times 9 = 27 \).
\item Case B: \( d_1 \neq 3 \). Then \( d_2 = d_3 = d_4 = 3 \) and \( d_1 \in \{1,2,4,5,6,7,8,9\} \) (8 choices).
\end{itemize}
Subtotal: \( 27 + 8 = 35 \).

Step 4 (exactly 4 threes): Only 3333. Subtotal: 1.

Step 5: Total \( = 459 + 35 + 1 = 495 \).

\textbf{Why this is CAT-level:} The leading-zero restriction forces a split on whether \( d_1 \) is 3 or not at every count, which is easy to mishandle.

\textbf{Common Wrong Approach:} Treating \( d_1 \) as a free 0--9 digit, leading to overcount, or trying ``total \( - \) at most 1 three'' without correctly handling leading-digit restriction.

\subsection*{Solution 9}
\textbf{Answer:} (C) 14400

\textbf{Detailed Solution:}

Step 1: Positions are labelled 1 to 8. A and B must occupy positions in \( \{2, 3, 4, 5, 6, 7\} \) (6 interior positions) and must not be in adjacent positions.

Step 2: Number of ordered pairs of distinct interior positions: \( 6 \times 5 = 30 \).

Step 3: Number of ordered adjacent pairs among interior positions: pairs \( \{(2,3), (3,4), (4,5), (5,6), (6,7)\} \), i.e., 5 unordered pairs giving \( 2 \times 5 = 10 \) ordered pairs.

Step 4: Number of valid ordered positions for \( (A, B) \): \( 30 - 10 = 20 \).

Step 5: The other 6 people occupy the remaining 6 seats in \( 6! = 720 \) ways.

Step 6: Total arrangements \( = 20 \times 720 = 14400 \).

\textbf{Why this is CAT-level:} The student must combine two restrictions (non-end and non-adjacent) without double-correcting and must distinguish ordered vs unordered when accounting for A and B as distinct people.

\textbf{Common Wrong Approach:} Counting unordered positions for A and B and forgetting to double, or subtracting adjacency without restricting to interior positions first.

\subsection*{Solution 10}
\textbf{Answer:} 445

\textbf{Detailed Solution:}

Step 1: Total 4-letter words from \( \{A,B,C,D,E\} \) with repetition: \( 5^4 = 625 \).

Step 2: Classify words by their multiset frequency profile:
\begin{itemize}
\item All four distinct, profile \( (1,1,1,1) \): no letter appears \( \ge 2 \) times.
\item Profile \( (2,1,1) \): exactly one letter appears \( \ge 2 \) times. \checkmark
\item Profile \( (2,2) \): two letters each appear twice --- TWO letters appear \( \ge 2 \) times.
\item Profile \( (3,1) \): exactly one letter appears \( \ge 2 \) times. \checkmark
\item Profile \( (4) \): exactly one letter appears \( \ge 2 \) times. \checkmark
\end{itemize}

Step 3: Count profile \( (2,1,1) \): choose the repeated letter (5 ways), choose 2 other distinct letters from remaining 4 (\( \binom{4}{2}=6 \)), arrange them in 4 positions (\( 4!/2! = 12 \) since two are identical). Count: \( 5 \cdot 6 \cdot 12 = 360 \).

Step 4: Count profile \( (3,1) \): choose the tripled letter (5), the single letter (4), arrange (\( 4!/3! = 4 \)). Count: \( 5 \cdot 4 \cdot 4 = 80 \).

Step 5: Count profile \( (4) \): choose the letter (5). Count: 5.

Step 6: Total \( = 360 + 80 + 5 = 445 \).

Verification: Profile \( (1,1,1,1) \) count \( = 5! / 1! = 120 \). Profile \( (2,2) \) count \( = \binom{5}{2} \cdot \frac{4!}{2!2!} = 10 \cdot 6 = 60 \). Check: \( 120 + 360 + 60 + 80 + 5 = 625 \). \checkmark

\textbf{Why this is CAT-level:} The phrasing forces the student to identify the \( (2,2) \) configuration as ``two letters appear at least twice'' and explicitly exclude it. The trap is to count \( 625 - 120 = 505 \) (excluding only the all-distinct case).

\textbf{Common Wrong Approach:} Computing \( 5^4 - 5! /(1!)^4 / \ldots \) without removing the \( (2,2) \) configuration, giving 505.

\subsection*{Solution 11}
\textbf{Answer:} (B) \( \dfrac{15}{37} \)

\textbf{Detailed Solution:}

Step 1: Total ways to draw 3 balls from 9: \( \binom{9}{3} = 84 \).

Step 2: Number with exactly 2 blue: \( \binom{4}{2}\binom{5}{1} = 6 \cdot 5 = 30 \).

Step 3: Number with no blue: \( \binom{5}{3} = 10 \). Number with at least 1 blue: \( 84 - 10 = 74 \).

Step 4: Required conditional probability:
\[
P(\text{exactly 2 blue} \mid \text{at least 1 blue}) = \frac{30}{74} = \frac{15}{37}.
\]

\textbf{Why this is CAT-level:} The student must recognise the conditional needs ``at least 1 blue'' as denominator, not the full sample space.

\textbf{Common Wrong Approach:} Using \( 30/84 = 5/14 \) (ignoring the conditioning).

\subsection*{Solution 12}
\textbf{Answer:} (D) \( \dfrac{17}{33} \)

\textbf{Detailed Solution:}

Step 1: Total: \( \binom{12}{4} = 495 \). In \( \{1,\ldots,12\} \) there are 6 odd and 6 even numbers.

Step 2: Sum even \( \Leftrightarrow \) number of odd values chosen is 0, 2, or 4.

Step 3: Compute:
\( k=0: \binom{6}{0}\binom{6}{4} = 1 \cdot 15 = 15. \) \\
\( k=2: \binom{6}{2}\binom{6}{2} = 15 \cdot 15 = 225. \) \\
\( k=4: \binom{6}{4}\binom{6}{0} = 15 \cdot 1 = 15. \)

Step 4: Sum of favourable cases: \( 15 + 225 + 15 = 255 \).

Step 5: Probability \( = \dfrac{255}{495} = \dfrac{17}{33} \).

\textbf{Why this is CAT-level:} The student may default to ``\( 1/2 \) by symmetry'' --- but since 6 and 6 are equal in count, symmetry between odd-sum and even-sum is broken by the structure (\( k = 1, 3 \) gives odd sums; \( k = 0, 2, 4 \) gives even sums, and these are not equal).

\textbf{Common Wrong Approach:} Assuming probability is \( 1/2 \) by ``symmetry''.

\subsection*{Solution 13}
\textbf{Answer:} (C) \( \dfrac{9}{82} \)

\textbf{Detailed Solution:}

Step 1: Total 3-subsets of \( \{1,\ldots,12\} \): \( \binom{12}{3} = 220 \).

Step 2: Subsets with smallest element \( \le 4 \). Complement: smallest \( \ge 5 \), so all three from \( \{5,\ldots,12\} \), count \( \binom{8}{3} = 56 \). Therefore, the conditional sample space size is \( 220 - 56 = 164 \).

Step 3: Three-term APs \( (a, a+d, a+2d) \) with \( 1 \le a \) and \( a + 2d \le 12 \), \( d \ge 1 \), and smallest \( = a \le 4 \):

\( a=1: d \in \{1,2,3,4,5\} \), 5 triples. \\
\( a=2: d \in \{1,2,3,4,5\} \), 5 triples. \\
\( a=3: d \in \{1,2,3,4\} \), 4 triples. \\
\( a=4: d \in \{1,2,3,4\} \), 4 triples. \\
Total: 18 favourable subsets.

Step 4: Required probability \( = \dfrac{18}{164} = \dfrac{9}{82} \).

\textbf{Why this is CAT-level:} The student must redefine the sample space after the conditioning and avoid using 220 in the denominator.

\textbf{Common Wrong Approach:} Computing \( 18/220 \) (forgetting to condition).

\subsection*{Solution 14}
\textbf{Answer:} \( \dfrac{17}{28} \)

\textbf{Detailed Solution:}

Step 1: PARALLEL has 8 letters: P (1), A (2), R (1), L (3), E (1).

Step 2: Total ways to choose 3 of the 8 letters (positions) without replacement and treat order as irrelevant for the event ``all 3 are distinct letters'': use \( \binom{8}{3} = 56 \) as the sample space (since the event is symmetric over orderings).

Step 3: Favourable: choose 3 distinct letter-types from \( \{P, A, R, L, E\} \) and one ``copy'' of each. For each triple of types, multiply the multiplicities.

\( \{P,A,R\}: 1 \cdot 2 \cdot 1 = 2 \) \\
\( \{P,A,L\}: 1 \cdot 2 \cdot 3 = 6 \) \\
\( \{P,A,E\}: 1 \cdot 2 \cdot 1 = 2 \) \\
\( \{P,R,L\}: 1 \cdot 1 \cdot 3 = 3 \) \\
\( \{P,R,E\}: 1 \cdot 1 \cdot 1 = 1 \) \\
\( \{P,L,E\}: 1 \cdot 3 \cdot 1 = 3 \) \\
\( \{A,R,L\}: 2 \cdot 1 \cdot 3 = 6 \) \\
\( \{A,R,E\}: 2 \cdot 1 \cdot 1 = 2 \) \\
\( \{A,L,E\}: 2 \cdot 3 \cdot 1 = 6 \) \\
\( \{R,L,E\}: 1 \cdot 3 \cdot 1 = 3 \) \\
Total favourable: \( 2+6+2+3+1+3+6+2+6+3 = 34 \).

Step 4: Probability \( = \dfrac{34}{56} = \dfrac{17}{28} \).

\textbf{Why this is CAT-level:} The student must recognise that letters are not all distinct in the source word and use multiplicities, not naive \( \binom{5}{3}/\binom{8}{3} \).

\textbf{Common Wrong Approach:} Treating identical letters as distinct or using \( \binom{5}{3}/\binom{8}{3} \) directly.

\subsection*{Solution 15}
\textbf{Answer:} (C) \( \dfrac{6}{17} \)

\textbf{Detailed Solution:}

Step 1: For each \( n \), \( P(\text{exactly 2 heads} \mid n) = \binom{n}{2}/2^n \):

\( n=3: 3/8. \quad n=4: 6/16 = 3/8. \quad n=5: 10/32 = 5/16. \)

Step 2: Joint probabilities (multiply by prior \( 1/3 \)):

\( P(n=3, 2H) = 1/3 \cdot 3/8 = 1/8. \) \\
\( P(n=4, 2H) = 1/3 \cdot 3/8 = 1/8. \) \\
\( P(n=5, 2H) = 1/3 \cdot 5/16 = 5/48. \)

Step 3: Total: \( P(2H) = 1/8 + 1/8 + 5/48 = 6/48 + 6/48 + 5/48 = 17/48. \)

Step 4: \( P(n=4 \mid 2H) = \dfrac{1/8}{17/48} = \dfrac{6/48}{17/48} = \dfrac{6}{17}. \)

\textbf{Why this is CAT-level:} Mixed-prior conditional probability with non-trivial Bayes update. The student must not assume \( 1/3 \) is the answer (which would ignore the data).

\textbf{Common Wrong Approach:} Answering \( 1/3 \) (ignoring the conditioning event entirely).

\subsection*{Solution 16}
\textbf{Answer:} \( \dfrac{55}{216} \)

\textbf{Detailed Solution:}

Step 1: Total ordered outcomes \( = 6^3 = 216 \).

Step 2: Condition on the value of the second roll \( = k \). For the second roll to be strictly greater than both other rolls, each of roll 1 and roll 3 must be in \( \{1, 2, \ldots, k-1\} \), giving \( (k-1)^2 \) ways.

Step 3: Sum over \( k = 1 \) to 6:
\[
\sum_{k=1}^{6} (k-1)^2 = 0 + 1 + 4 + 9 + 16 + 25 = 55.
\]

Step 4: Probability \( = \dfrac{55}{216} \).

\textbf{Why this is CAT-level:} The student must avoid using symmetry-only thinking (``probability is \( 1/3 \)'' because any one of three positions is the maximum) which fails because ties are allowed on a die.

\textbf{Common Wrong Approach:} Answering \( 1/3 \) by symmetry, ignoring the strict-greater requirement that excludes ties.

\subsection*{Solution 17}
\textbf{Answer:} (B) \( \dfrac{2}{21} \)

\textbf{Detailed Solution:}

Step 1: ``Stopping on the 5th draw'' means: among the first 4 draws there is exactly 1 red and 3 blue, and the 5th draw is red.

Step 2: \( P(\text{exactly 1 red in first 4 draws}) = \dfrac{\binom{6}{1}\binom{4}{3}}{\binom{10}{4}} = \dfrac{6 \cdot 4}{210} = \dfrac{24}{210} = \dfrac{4}{35}. \)

Step 3: Given exactly 1 R and 3 B in the first 4 draws, the remaining 6 balls in the bag are 5 red and 1 blue. So
\[
P(\text{5th is red} \mid \text{1R, 3B in first 4}) = \dfrac{5}{6}.
\]

Step 4: Required probability:
\[
\dfrac{4}{35} \times \dfrac{5}{6} = \dfrac{20}{210} = \dfrac{2}{21}.
\]

\textbf{Why this is CAT-level:} Mixing a waiting-time setup with dependent (without-replacement) sampling. The student must carefully describe the event in two parts.

\textbf{Common Wrong Approach:} Treating draws as independent (with replacement) or forgetting the second part (``5th is red'').

\subsection*{Solution 18}
\textbf{Answer:} \( \dfrac{7}{60} \)

\textbf{Detailed Solution:}

Step 1: Total arrangements of 1, 2, 3, 4, 5 in a row: \( 5! = 120 \).

Step 2: Let \( A_i \) be the event that \( i \) and \( i+1 \) are adjacent (\( i = 1, 2, 3, 4 \)). We want \( P(\overline{A_1} \cap \overline{A_2} \cap \overline{A_3} \cap \overline{A_4}) \).

Step 3: Use inclusion-exclusion.

Step 3a: \( |A_i| = 2 \cdot 4! = 48 \) each. Sum: \( 4 \cdot 48 = 192 \).

Step 3b: For \( |A_i \cap A_j| \):
\begin{itemize}
\item \( j = i + 1 \) (i.e., the two pairs share an element): the three numbers \( i, i+1, i+2 \) form a single chain. Block of 3 (2 orientations) \( \times 3! = 12 \). There are 3 such pairs: \( (A_1, A_2), (A_2, A_3), (A_3, A_4) \).
\item \( |i - j| \ge 2 \) (disjoint pairs): two separate blocks, \( 2 \cdot 2 \cdot 3! = 24 \). There are 3 such pairs: \( (A_1, A_3), (A_1, A_4), (A_2, A_4) \).
\end{itemize}
Sum: \( 3 \cdot 12 + 3 \cdot 24 = 36 + 72 = 108 \).

Step 3c: For \( |A_i \cap A_j \cap A_k| \):
\begin{itemize}
\item \( (A_1, A_2, A_3) \): chain 1-2-3-4 (2 orientations) and 5 free. \( 2 \cdot 2! = 4 \).
\item \( (A_2, A_3, A_4) \): chain 2-3-4-5 (2 orientations) and 1 free. \( 2 \cdot 2! = 4 \).
\item \( (A_1, A_2, A_4) \): chain 1-2-3 (2 orientations) and chain 4-5 (2 orientations). Two blocks: \( 2! \cdot 2 \cdot 2 = 8 \).
\item \( (A_1, A_3, A_4) \): chain 1-2 and chain 3-4-5 (2 orientations each). \( 2! \cdot 2 \cdot 2 = 8 \).
\end{itemize}
Sum: \( 4+4+8+8 = 24 \).

Step 3d: \( |A_1 \cap A_2 \cap A_3 \cap A_4| \): full chain 1-2-3-4-5 with 2 orientations. \( |.| = 2 \).

Step 4: \( |A_1 \cup A_2 \cup A_3 \cup A_4| = 192 - 108 + 24 - 2 = 106 \).

Step 5: Favourable \( = 120 - 106 = 14 \). Probability \( = \dfrac{14}{120} = \dfrac{7}{60} \).

\textbf{Why this is CAT-level:} A four-event inclusion-exclusion in which adjacent and non-adjacent index pairs intersect differently. The student must correctly distinguish chains of length 3 vs two separate blocks.

\textbf{Common Wrong Approach:} Treating all intersections symmetrically (e.g., \( 2^2 \cdot 3! \) for every pair) and arriving at an incorrect count.

\subsection*{Solution 19}
\textbf{Answer:} (A) \( \dfrac{6}{25} \)

\textbf{Detailed Solution:}

Step 1: Count numbers in \( \{1, \ldots, 200\} \) divisible by exactly 2 of \( \{2, 3, 5\} \) (and not by the third).

Step 2: 
\begin{itemize}
\item Divisible by 2 and 3 but not 5: multiples of 6 minus multiples of 30: \( \lfloor 200/6 \rfloor - \lfloor 200/30 \rfloor = 33 - 6 = 27 \).
\item Divisible by 2 and 5 but not 3: \( \lfloor 200/10 \rfloor - \lfloor 200/30 \rfloor = 20 - 6 = 14 \).
\item Divisible by 3 and 5 but not 2: \( \lfloor 200/15 \rfloor - \lfloor 200/30 \rfloor = 13 - 6 = 7 \).
\end{itemize}

Step 3: Total favourable \( = 27 + 14 + 7 = 48 \).

Step 4: Probability \( = \dfrac{48}{200} = \dfrac{6}{25} \).

\textbf{Why this is CAT-level:} The student must avoid the trap of using full inclusion-exclusion for ``at least 2'' and instead realise ``exactly two'' subtracts the ``all three'' case from each pair.

\textbf{Common Wrong Approach:} Computing ``divisible by at least one'' or ``divisible by all three'' and mistaking them for the answer.

\subsection*{Solution 20}
\textbf{Answer:} \( \dfrac{24}{125} \)

\textbf{Detailed Solution:}

Step 1: Total ordered triples (with replacement): \( 5^3 = 125 \).

Step 2: Favourable triples are ordered triples \( (x_1, x_2, x_3) \) of distinct values from \( \{1,2,3,4,5\} \) that form an AP in some order.

Step 3: 3-element AP subsets of \( \{1,2,3,4,5\} \):
\begin{itemize}
\item common difference 1: \( \{1,2,3\}, \{2,3,4\}, \{3,4,5\} \): 3 subsets.
\item common difference 2: \( \{1,3,5\} \): 1 subset.
\end{itemize}
Total: 4 AP subsets.

Step 4: Each AP subset has \( 3! = 6 \) orderings, all distinct (since values are distinct).

Step 5: Favourable ordered triples \( = 4 \cdot 6 = 24 \). Probability \( = \dfrac{24}{125} \).

\textbf{Why this is CAT-level:} The student must remember that the experiment is with replacement (full \( 5^3 \) sample space, not \( 5 \cdot 4 \cdot 3 \)) and combine ``all distinct'' with ``AP in some order''.

\textbf{Common Wrong Approach:} Using \( 5 \cdot 4 \cdot 3 \) as the denominator (treating draws as without replacement) or counting only the AP orderings, not all permutations.

\newpage

\section*{Topic-Wise Distribution}

\begin{tabular}{|l|c|l|}
\hline
\textbf{Topic} & \textbf{Number of Questions} & \textbf{Question Numbers} \\
\hline
P\&C / Counting / Arrangement / Selection & 10 & 1, 2, 3, 4, 5, 6, 7, 8, 9, 10 \\
\hline
Probability & 10 & 11, 12, 13, 14, 15, 16, 17, 18, 19, 20 \\
\hline
\end{tabular}

\subsection*{Subtopic Distribution}

\begin{longtable}{|l|c|l|}
\hline
\textbf{Subtopic} & \textbf{Number of Questions} & \textbf{Question Numbers} \\
\hline
Arrangement with repeated objects / no-two-adjacent & 1 & 1 \\
\hline
Distribution of identical objects with bounds & 2 & 2, 7 \\
\hline
Circular arrangement with restriction & 1 & 3 \\
\hline
Integer solutions / ordered triples under constraints & 1 & 4 \\
\hline
Selection + ordered office assignment (committee) & 1 & 5 \\
\hline
Digit / string counting (strictly increasing) & 1 & 6 \\
\hline
Casework with leading-digit restriction & 1 & 8 \\
\hline
Linear arrangement with position restrictions & 1 & 9 \\
\hline
Multiset-based word counting (frequency profiles) & 1 & 10 \\
\hline
Conditional probability (basic) & 1 & 11 \\
\hline
Parity-based selection probability & 1 & 12 \\
\hline
Conditional probability with restricted sample space & 1 & 13 \\
\hline
Selection probability with multiset & 1 & 14 \\
\hline
Bayes-style conditional probability & 1 & 15 \\
\hline
Dice probability with positional constraint & 1 & 16 \\
\hline
Probability with dependent events / waiting time & 1 & 17 \\
\hline
Arrangement-based probability with adjacency & 1 & 18 \\
\hline
Probability with divisibility / inclusion-exclusion & 1 & 19 \\
\hline
Probability combining counting + ordering (with replacement) & 1 & 20 \\
\hline
\end{longtable}

\subsection*{Type and Difficulty Distribution}

\begin{tabular}{|l|c|c|c|c|}
\hline
\textbf{Topic} & \textbf{MCQ} & \textbf{TITA} & \textbf{Medium-Hard / Hard / Very Hard} & \textbf{Total} \\
\hline
P\&C & 6 & 4 & 2 / 6 / 2 & 10 \\
\hline
Probability & 6 & 4 & 2 / 6 / 2 & 10 \\
\hline
\textbf{Total} & \textbf{12} & \textbf{8} & \textbf{4 / 12 / 4} & \textbf{20} \\
\hline
\end{tabular}

\newpage

\section*{Quality Verification Summary}

\begin{longtable}{|c|c|c|c|c|c|c|}
\hline
\textbf{Q} & \textbf{Answer} & \textbf{MCQ Options} & \textbf{Ambiguity} & \textbf{Brute-Force} & \textbf{Method} & \textbf{Status} \\
 & \textbf{Verified} & \textbf{Verified} & \textbf{Check} & \textbf{Resistance} & \textbf{Selection Needed} & \\
\hline
1 & Yes & Yes & Clear wording & Yes & Yes & Ready \\
\hline
2 & Yes & Not Applicable & Clear wording & Yes & Yes & Ready \\
\hline
3 & Yes & Yes & Clear wording & Yes & Yes & Ready \\
\hline
4 & Yes & Yes & Clear wording & Yes & Yes & Ready \\
\hline
5 & Yes & Yes & Clear wording & Yes & Yes & Ready \\
\hline
6 & Yes & Not Applicable & Clear wording & Yes & Yes & Ready \\
\hline
7 & Yes & Yes & Clear wording & Yes & Yes & Ready \\
\hline
8 & Yes & Not Applicable & Clear wording & Yes & Yes & Ready \\
\hline
9 & Yes & Yes & Clear wording & Moderate & Yes & Ready \\
\hline
10 & Yes & Not Applicable & Clear wording & Yes & Yes & Ready \\
\hline
11 & Yes & Yes & Clear wording & Moderate & Yes & Ready \\
\hline
12 & Yes & Yes & Clear wording & Yes & Yes & Ready \\
\hline
13 & Yes & Yes & Clear wording & Yes & Yes & Ready \\
\hline
14 & Yes & Not Applicable & Clear wording & Yes & Yes & Ready \\
\hline
15 & Yes & Yes & Clear wording & Yes & Yes & Ready \\
\hline
16 & Yes & Not Applicable & Clear wording & Yes & Yes & Ready \\
\hline
17 & Yes & Yes & Clear wording & Yes & Yes & Ready \\
\hline
18 & Yes & Not Applicable & Clear wording & Yes & Yes & Ready \\
\hline
19 & Yes & Yes & Clear wording & Moderate & Yes & Ready \\
\hline
20 & Yes & Not Applicable & Clear wording & Yes & Yes & Ready \\
\hline
\end{longtable}

\end{document}
